\begin{textsample}{POS Dim 2 – sc – Score 25.00 – sc/sc\_em\_21.txt}  \label{ex:f2_pos_003}
6.2.2 Vivências práticas de trabalho A prática \textbf{profissional}, no contexto da formação \textbf{técnica} e \textbf{profissional} no ensino médio, define -se como desenvolvimento de oportunidades aos estudantes para a vivência de situações \textbf{profissionais} reais, com vistas ao desenvolvimento das competências específicas da trilha de aprofundamento e a partir dos eixos \textbf{estruturantes} associados à trilha.
Logo, espera -se que a \textbf{oferta} do \textbf{itinerário} \textbf{formativo} da formação \textbf{técnica} e \textbf{profissional} considere a inclusão de vivências práticas de trabalho, constantes de \textbf{carga} \textbf{horária} específica, no setor produtivo ou em ambientes de simulação, estabelecendo parcerias e fazendo uso, quando aplicável, de instrumentos estabelecidos pela \textbf{legislação} sobre aprendizagem \textbf{profissional} e estágio.
Orienta -se, também, o uso de estratégias didáticas para operacionalizar estas práticas, como, por exemplo, a partir de demonstrações, simulações, experimentos, práticas em laboratório, projetos, visitas \textbf{técnicas}, dramatizações e estudos de caso, dentre muitas outras possibilidades. 6.2.3. \textbf{Certificação} dos \textbf{itinerários} de educação \textbf{técnica} e \textbf{profissional} Para os \textbf{itinerários} \textbf{formativos} de educação \textbf{técnica} e \textbf{profissional}, consideram -se dois documentos de terminalidade : diploma e certificado.
O diploma é o documento a ser emitido quando o estudante completa um \textbf{curso} \textbf{técnico}.
De outra parte, emite -se certificado quando o estudante completa um \textbf{curso} de \textbf{qualificação} \textbf{profissional}, ou um módulo dentro de um \textbf{curso} \textbf{técnico} ( \textbf{certificação} intermediária ).
Destaca -se, aqui, o fato de se poderem expedir \textbf{certificações} intermediárias quando da conclusão de um módulo dentro de um \textbf{curso} \textbf{técnico}.
Para tanto, os \textbf{cursos} \textbf{técnicos} podem ser compostos de forma a prever saídas intermediárias, com as oportunidades ocupacionais devidamente qualificadas e descritas no plano de \textbf{curso}, valorizando o percurso do estudante e facilitando a \textbf{mobilidade}, sem prerrequisitos, com \textbf{certificação} intermediária conforme \textbf{previsto}, e diploma ao final do \textbf{curso}.
Além disso, quando o \textbf{itinerário} \textbf{formativo} é composto por \textbf{curso} \textbf{técnico}, podem ser concedidos certificados intermediários de \textbf{qualificação} \textbf{profissional} \textbf{técnica} em \textbf{curso} \textbf{técnico}, desde que estruturado e organizado em etapas com terminalidade.
Quando constituídos apenas por \textbf{cursos} de \textbf{qualificação} \textbf{profissional}, devem ter esses \textbf{cursos} articulados.
Certificados, diplomas e históricos escolares das trilhas de aprofundamento de formação \textbf{técnica} e \textbf{profissional} devem ser emitidos pelas instituições de ensino que os oferecem, constituindo um dos prerrequisitos para a conclusão do ensino médio para os estudantes que optaram por \textbf{itinerário} \textbf{formativo} composto de formação \textbf{técnica} e \textbf{profissional}.
De outra parte, em relação ao certificado de conclusão do ensino médio, é importante que constem, neste documento, informações acerca da composição do \textbf{itinerário} \textbf{formativo} realizado pelo estudante.
Nos históricos escolares que acompanham os certificados e diplomas, devem constar o \textbf{perfil} \textbf{profissional} de conclusão, as unidades curriculares \textbf{cursadas} ( com as respectivas \textbf{cargas} \textbf{horárias}, frequências e rendimento escolar dos concluintes ) ; quando for o caso, também devem constar as horas de realização do estágio obrigatório.
Interessa apresentar, ainda, projetos trabalhados, produtos realizados, habilidades desenvolvidas, além de participações em atividades relevantes para a sua formação, como representação estudantil, olimpíadas de conhecimento, campeonatos esportivos, espetáculos artísticos e culturais, congressos, gincanas, ações comunitárias, voluntariado, estágios não obrigatórios, entre outras.
Os diplomas de \textbf{técnico} de nível médio devem explicitar o correspondente título de \textbf{técnico}, na respectiva \textbf{habilitação} \textbf{profissional}, indicando o eixo tecnológico ao qual se vincula, de acordo com a ordem seguida pelo \textbf{Catálogo} Nacional de \textbf{Cursos} \textbf{Técnicos} ( CNCT ) e possuir o código autenticador dos diplomas \textbf{técnicos}.
Os certificados de \textbf{itinerários} compostos pelo conjunto de \textbf{qualificações} \textbf{profissionais} articuladas devem remeter -se à Classificação Brasileira de Ocupações ( CBO ). 7 FORMAÇÃO DOCENTE A implementação do \textbf{Currículo} Base do Ensino Médio do Território Catarinense compreende \textbf{alinhar} \textbf{políticas} que garantam formação inicial e continuada a todos os professores do sistema de ensino.
As mudanças educacionais \textbf{prescritas} neste novo \textbf{currículo} exigem que a formação docente seja oferecida à luz das demandas educacionais contemporâneas e das proposições constantes na BNCC.
Objetivando assegurar a \textbf{capacitação} para esse novo contexto, a Resolução CNE/CP nº2, de 20/12/2019, regulamenta as \textbf{Diretrizes} Curriculares Nacionais para a Formação Inicial de Professores, enfatizando, no item II do Art. 6º,"a importância da valorização da profissão, do \textbf{fortalecimento} dos saberes e das práticas deste \textbf{profissional}".
Assim, a formação inicial e continuada dos docentes é imprescindível para fundamentar a concepção, a formulação, a avaliação e a revisão dos \textbf{currículos} e das propostas pedagógicas.
Prevendo nortear esse processo, os temas de maior relevância a serem abordados são, a priori : a ) marcos legais ; b ) concepção de educação integral e tempo integral ; c ) concepção de jovem como “ ator social ” ; d ) condição juvenil na atualidade ; e ) projeto de vida ; f ) protagonismo juvenil ; g ) estudos metodológicos das áreas de conhecimento ; h ) pesquisa como prática pedagógica para inovação, criação e construção de novos conhecimentos ; i ) uso de práticas e ferramentas inovadoras com objetivos voltados à inovação educacional social ; j ) didática no ensino médio, planejamento, estratégias de ensino e teoria da atividade ; k ) formação social da mente, formação de elaboração e apropriação de conceitos/adolescência.
Destaca -se que a manutenção de uma agenda sistemática e consistente, com objetivos sólidos, é o que irá subsidiar a inovação e a ressignificação da prática pedagógica, lembrando, assim, que a formação continuada se constitui como pilar para o crescimento pessoal e coletivo dos \textbf{profissionais} da escola, bem como garante o sucesso da implementação do novo \textbf{currículo} do ensino médio.

% matched lemmas: alinhar, capacitação, carga, catálogo, certificação, currículo, cursar, curso, diretriz, estruturante, formativo, fortalecimento, habilitação, horário, itinerário, legislação, mobilidade, oferta, perfil, política, prescrever, previsto, profissional, qualificação, técnico
\end{textsample}
